% Generated draft: requires mathematical and source review.
\documentclass[11pt]{article}
\usepackage[margin=25mm]{geometry}
\usepackage{fontspec}
\setmainfont{Latin Modern Roman}
\usepackage{amsmath,amssymb,mathtools,unicode-math}
\setmathfont{Latin Modern Math}
\usepackage{enumitem,needspace,xurl,hyperref}
\setlength{\parindent}{0pt}
\setlength{\parskip}{5pt}
\setlength{\emergencystretch}{3em}
\newcommand{\N}{\mathbb{N}}
\newcommand{\Z}{\mathbb{Z}}
\newcommand{\Q}{\mathbb{Q}}
\newcommand{\R}{\mathbb{R}}
\newcommand{\C}{\mathbb{C}}
\newcommand{\F}{\mathbb{F}}
\newcommand{\A}{\mathbb{A}}
\newcommand{\PP}{\mathbb{P}}
\newcommand{\E}{\mathbb{E}}
\newcommand{\eps}{\varepsilon}
\newcommand{\dd}{\,\mathrm{d}}
\begin{document}
\section*{UnsolvedMath ID: 30000389}
\textbf{Source problem code:} catalogue entry
\subsection{Definitions and mathematical statement}
Let $G=(V,E)$ be a finite simple undirected graph, let $n=|V|$, and assume $G$ is triangle-free, meaning that no three distinct vertices span all three possible edges. Define its edge bipartization number by
\[
 b(G)=\min\{|F|:F\subseteq E,\ (V,E\setminus F)\text{ is bipartite}\}.
\]
A graph is bipartite if its vertices can be partitioned into two sets with no edges within either set. The conjecture attributed in the supplied assessment to Erd\H{o}s asks whether
\[
 \forall n\in\N\ \forall G\text{ triangle-free with }|V(G)|=n:
 \qquad b(G)\le\frac{n^2}{25}.
\]
Since $b(G)$ is an integer, the exact bound is equivalently $b(G)\le\lfloor n^2/25\rfloor$. Only edges are deleted; all $n$ vertices remain.

Equivalently, writing $\operatorname{MaxCut}(G)$ for the largest number of edges crossing a partition $V=A\sqcup B$, the question is
\[
 \operatorname{MaxCut}(G)\ge |E(G)|-\frac{n^2}{25}.
\]
There is no asymptotic error term or restriction to sufficiently large $n$. The finite simple unweighted graph convention is the standard interpretation of the source; loops, parallel edges, and weighted variants are not supplied.
\subsection{Short English statement}
Can every finite triangle-free graph be made bipartite by deleting at most $n^2/25$ edges, with the bound exact for every number of vertices? Equivalently, its largest cut must leave at most that many edges uncut.
\subsection{Sources}
\begin{itemize}
\item Local UnsolvedMath registry, permanent ID 30000389, code catalogue entry, ``Bipartizing Triangle-Free Graphs with Few Edge Deletions.'' Proposed year 2006; proposer field empty. The literature assessment calls this the Erd\H{o}s bipartization conjecture.
\item \emph{Combinatorics} (2006), Oberwolfach Reports, pp.\ 5--72, DOI 10.4171/owr/2006/01; \url{https://doi.org/10.4171/owr/2006/01}.
\item Supplied reference ``Max Cuts in triangle-free graphs'': \url{https://arxiv.org/abs/2103.14179}.
\item Supplied reference ``Erdos n squared over 25 max-cut conjecture'': \url{https://arxiv.org/abs/2606.28041}.
\item Statement status corrected\_verified; the supplied verification says the official report was checked and wording made self-contained with normalized notation.
\item Literature triage checked 2026-08-21: partially solved. It describes the full conjecture as open, reports density-tail results and improved global bounds, and says 2026 work verifies additional finite sizes. No numerical ranges or theorem hypotheses are supplied. Created 2026-08-25; updated 2026-09-16. Time-sensitive claims require independent verification.
\item Registry curation by Ulam AI / UnsolvedMath Contributors, CC BY 4.0; underlying report and literature provenance remain as supplied. Unreviewed statement-only draft.
\end{itemize}
\end{document}
